\documentclass[10pt]{article}\usepackage[margin=0.8in]{geometry}\usepackage{amssymb,longtable,array,url,hyperref}\usepackage[T1]{fontenc}\usepackage{lmodern}\sloppy\usepackage{newunicodechar}\newunicodechar{η}{\ensuremath{\eta}}\newunicodechar{ω}{\ensuremath{\omega}}\newunicodechar{⁺}{\ensuremath{^{+}}}\newunicodechar{⁻}{\ensuremath{^{-}}}\newunicodechar{∏}{\ensuremath{\prod}}\newunicodechar{⁸}{\ensuremath{^{8}}}\newunicodechar{⁴}{\ensuremath{^{4}}}\newunicodechar{⁶}{\ensuremath{^{6}}}\newunicodechar{θ}{\ensuremath{\theta}}\newunicodechar{₇}{\ensuremath{_{7}}}\newunicodechar{₂}{\ensuremath{_{2}}}\newunicodechar{₁}{\ensuremath{_{1}}}\newunicodechar{⊕}{\ensuremath{\oplus}}
\begin{document}
\begin{center}{\Large\bfseries s-family: internal patterns, pre-registered observable matching, structural links}\end{center}

Prepared 20 Oct 2026 (CDT). Local only: nothing was posted, the repo was not edited and nobody was contacted. All computations use PARI (cypari2), numpy and scipy.

Files:
\begin{itemize}
\item  \texttt{partA\_table.csv} (s = 5..500), \texttt{partA\_analysis.out}, \texttt{partA\_correlations.csv}.
\item  \texttt{PREREG.md}, with sha256 and timestamp in \texttt{prereg.sha256}. It was written before \texttt{partB\_match.py} was run.
\item  \texttt{partB\_results.json}, \texttt{partB\_matches.csv}, \texttt{partB\_posthoc.json}.
\end{itemize}

Part A is labelled Category A where proved, Category B otherwise. Part C is Category B, speculative.

\section*{Part A: internal patterns (s = 5..500)}

The table was built with the manuscript's definitions (s01\_intro.tex:131-134). The cusp dimension uses d\_cusp = dim V - d\_Eis, where:
\begin{itemize}
\item  dim V is the orbit bound s(s-1) for odd s and (s/2)(s/2+1) for even s, reduced by m(m-1) when s = 2m$^2$ (s12\_dims.tex:8-15);
\item  d\_Eis is the manuscript formula (s12\_dims.tex:16-19);
\item  the manuscript value dim V = 2406 is used at s = 98.
\end{itemize}

\textbf{Check:} these formulas reproduce all 48 tabulated (dim V, d\_Eis, d\_cusp) rows exactly (s = 5..51 and 72). For s > 51 they are extrapolations, so Category B. The manuscript proves exactness only for s $\leq$ 51 and certifies d\_Eis only for the listed local factors.

\subsection*{Proved (Category A)}

\begin{itemize}
\item  \textbf{When $\varepsilon$\_s is not the fundamental unit.} Write $\varepsilon$\_s = η\^{}k with η the fundamental unit and t its trace. Then:
\begin{itemize}
\item  for N(η) = +1, $\varepsilon$\_s = η\^{}k exactly when s - 1 = T\_k(t/2) (T\_k the Chebyshev polynomial);
\item  for N(η) = -1, the analogous condition holds with the norm -1 trace.
\end{itemize}
\end{itemize}

   This is standard: if η = (t+u$\surd$d)/2, then η\^{}k + η\^{}{-k} = 2T\_k(t/2). The 34 cases with k > 1 for s $\leq$ 500 are listed in \texttt{partA\_analysis.out}, and every one satisfies the condition with the computed integer trace. The cases fall into these families:
\begin{itemize}
\begin{itemize}
\item  k = 2, N = +1: s = 2m$^2$, for example 8, 32, 50, 72, 128 (the manuscript's mechanism, s12\_dims.tex:14-15). At s = 18, 98 and 162 the power is 4 or 12, not 2.
\item  k = 2, N = -1: s = 2m$^2$ + 2, i.e. 20, 34, 52, 74, \ldots{}.
\item  k = 3: s - 1 = (t$^3$ - 3t)/2, i.e. s = 27, 56, 245, 352, 486.
\item  k = 5: s = 363 (t = 4).
\item  k = 6 or 12: s = 10, 100, 162.
\end{itemize}
\end{itemize}

   In total $\varepsilon$\_s is fundamental for 462 of 496 values of s.
\begin{itemize}
\item  \textbf{Genus theory.} 2\^{}{ω(d\_K)-1} divides h⁺(K) for every s, verified 100\%. It follows that h⁺ is divisible by 4 on the twin-prime subfamily (d\_K = 4(s-2)s, ω = 3), as observed for all 24 twin-prime s up to 500.
\item  \textbf{Closed form for odd s,} given the manuscript's Theorem dims (iii) and dim V = s(s-1): d\_cusp = s(s-1) - 3∏\_{p\^{}e‖n\_s}(p\^{}e+1)/2.
\end{itemize}

\subsection*{Empirical (Category B)}

\begin{itemize}
\item  \textbf{Zero cusp part:} d\_cusp = 0 only at s = 6, 8, 10, both in the tabulated range and in the extrapolation to 500.
\item  \textbf{Parity:} d\_cusp is even for every s $\leq$ 500. This is not proved here.
\item  \textbf{Growth:} d\_cusp/dim V tends upward. log(ratio) $\approx$ 0.24 log s + const for odd s (r = 0.55) and 0.33 log s + const for even s (r = 0.64). The median ratio for s > 200 is 0.59 (odd) and 0.61 (even). This is a trend, not a closed-form law.
\item  \textbf{Cohen--Lenstra:} among the 479 distinct fields, 53.9\% have trivial odd part of h. The Cohen--Lenstra heuristic for real quadratic fields gives about 75.4\% (95\% half-width 3.9\%), so these fields depart from it. That is expected rather than anomalous: d = m$^2$ - 1 (m = s - 1) is Richaud--Degert type, with an unusually small regulator ($\varepsilon$\_s $\approx$ 2$\surd$d), so class numbers run larger (Yokoi; Mollin, *Quadratics*, 1996). The heuristic assumes generic regulators.
\item  \textbf{Correlations:} 36 Spearman tests (raw, and with the log-s trend removed), with 2000-sample permutation p-values and BH correction; see \texttt{partA\_correlations.csv}. Every BH-significant correlation follows from known relations:
\begin{itemize}
\item  h⁺ against the ray-class number ($\rho$ = 0.91): the ray-class exact sequence.
\item  h⁺ against the conductor f ($\rho$ = -0.85 detrended): the class number formula h$\cdot$R \~{} $\surd$d\_K, with d\_K = 4n\_s/f$^2$.
\item  regulator against k ($\rho$ = -0.97 detrended): R = log $\varepsilon$\_s / k.
\item  correlations with s itself: Brauer--Siegel growth.
\end{itemize}
\end{itemize}

  No correlation was found that is unexplained.

\section*{Part B: pre-registered observable matching}

\textbf{Setup} (PREREG.md):
\begin{itemize}
\item  Quantities: Q1 = 2 log $\varepsilon$\_s, Q2 = I\_s, Q3 = h⁺/h, Q4 = d\_cusp/dim V (tabulated s), Q5 = $\surd$n\_s. That is 2032 values.
\item  Operations: x, 1/x, 2x, 3x, x/2 and x/3.
\item  16 cited targets.
\item  \textbf{N = 195 072 trials.}
\item  Criterion: |f(x) - y| $\leq$ max($\sigma$, 10⁻$^3$|y|).
\end{itemize}

\textbf{Deviation from the pre-registration, disclosed:} T15's $\sigma$ was recomputed from the cited inputs as 54.6 ($\times$10⁻⁸), not the 58.4 typed in PREREG. That is an arithmetic slip in the pre-registration and does not change any conclusion.

\subsection*{Pre-registered result}

\begin{itemize}
\item  \textbf{Observed M = 10 633 matches.}
\item  Null means (10⁴ trials each):
\begin{itemize}
\item  N1 log-uniform: 11 903 $\pm$ 45, p = 1.0.
\item  N2 random-field units: 8822 $\pm$ 44, p $\approx$ 1e-4.
\item  N3 random integers: 10 505 $\pm$ 23, p $\approx$ 1e-4.
\end{itemize}
\item  \textbf{The global count is dominated by two pre-registered targets that turned out to be degenerate:}
\begin{itemize}
\item  T15 (g-2 deviation 33 $\pm$ 55): any value in [-22, 88] matches, giving 10 067 of the matches.
\item  T16 (m\_W CDF/PDG = 1.0008, with tolerance 10⁻$^3$): any value equal to 1 matches, so the 496 values of Q3 = 1 all match.
\end{itemize}
\item  The low p-values under N2/N3 come from differences between those nulls and the s-family value distribution inside the wide T15 window. They do not reflect precise matches. The p-values disagree across nulls (1.0 against 1e-4), which shows the global statistic is null-model-dependent and uninformative.
\end{itemize}

\subsection*{Post-hoc (labelled; \texttt{partB\_posthoc.py})}

Excluding T15 and T16:

{\small\begin{longtable}{p{0.04\linewidth}p{0.36\linewidth}p{0.44\linewidth}p{0.08\linewidth}}\hline
\textbf{} & \textbf{Count}\\\hline\endhead
Observed & 67\\\hline
N1 & 172.8 $\pm$ 13.1 (p = 1.0)\\\hline
N2 & 90.1 $\pm$ 9.5 (p = 0.99)\\\hline
N3 & 145.3 $\pm$ 11.6 (p = 1.0)\\\hline
N4 (jitter: each true value $\times$ e\^{}{N(0,0.02)}, closest-matched distribution) & 69.2 $\pm$ 6.1 (p = 0.67)\\\hline
\end{longtable}}

\textbf{The s-family produces no more matches than distribution-matched random numbers.}

Per target:
\begin{itemize}
\item  No target is significant under N4.
\item  T1 ($\alpha$⁻¹) had p = 0.04 under N1 only. With 16 $\times$ 4 tests that is not significant, and under N2, N3 and N4 the p-values are 0.22--0.48.
\item  \textbf{None of the 67 matches is within 1$\sigma$ for any precision target (T1--T5).} They pass only through the 10⁻$^3$ relative tolerance.
\end{itemize}

Examples:
\begin{itemize}
\item  $\surd$n\_s at s = 138 is 136.996, against $\alpha$⁻¹ = 137.036: a relative difference of 2.9e-4, i.e. 1.9$\times$10⁶ $\sigma$. This is a trivial near-integer, since $\surd$(s(s-2)) $\approx$ s - 1.
\item  1/I\_s at s = 470 is 137.059 (1.7e-4, about 10⁶ $\sigma$).
\item  1/I\_s at s = 37 is 16.830, against m\_$\tau$/m\_$\mu$ = 16.817 (0.08\%, 12$\sigma$).
\item  (2/2)$\cdot$d\_cusp/dim V at s = 41, giving 0.2311, against sin$^2$θ\_W = 0.23129 (0.08\%, 5$\sigma$).
\end{itemize}

\textbf{Verdict for Part B: chance.}

The value s = 7 itself matches nothing: 2 log $\varepsilon$₇ = 4.9558, I₇ = 0.20942, $\surd$35 = 5.916, d\_cusp/dim V = 1/7.

\section*{Part C: structural correspondences (Category B, speculative)}

\begin{itemize}
\item  \textbf{Weil representation as a quantum cat map (most promising).}
\begin{itemize}
\item  $\gamma$\_s = [[-1,-s],[2,2s-1]] is a hyperbolic SL₂(Z) element with entropy log $\varepsilon$\_s, an Anosov cat map.
\item  Its quantization is given by the Weil (metaplectic) representation: Hannay \& Berry, Physica D 1 (1980) 267; Kurlberg \& Rudnick, Duke Math J 103 (2000) 47 (Hecke eigenbases).
\item  The Bridge C label space carries exactly a Weil representation of the discriminant form (s03\_main.tex:77-87; s12\_dims.tex proof notes).
\item  \textbf{Next computation:} the spectrum of $\rho$\_s($\gamma$\_s) on W\_P for s = 5..60. Compare the eigenphase statistics and the period of $\rho$\_s($\gamma$\_s) mod the label level with Kurlberg--Rudnick predictions. Test whether the cusp sector V\_cusp coincides with an anti-invariant or Hecke-eigen subspace.
\end{itemize}
\item  \textbf{Indefinite theta series and wall-crossing / two-centred black-hole bound states.}
\begin{itemize}
\item  Signature-(1,1) Zwegers theta functions count two-centre D4-D2-D0 bound states: Manschot, arXiv:0906.1767 and arXiv:1009.1775; Alexandrov--Banerjee--Manschot--Pioline, arXiv:1606.05495 (generalized error functions).
\item  The mock modularity of BH microstate generating functions: Dabholkar--Murthy--Zagier, arXiv:1208.4074.
\item  \textbf{Next computation:} identify whether a lattice ⟨2⟩⊕⟨-2(s-2)s⟩-type form, with the cone walls c₁ = (-1,1) and c₂ = (-4,1), arises as the charge lattice of a two-centre configuration (some CY$_3$ intersection form, with D = s(s-2)). If it does, compare H\_s with the corresponding index jump across the wall.
\end{itemize}
\item  \textbf{Quantum modularity and 3-manifold invariants.}
\begin{itemize}
\item  The radial and rational values $\psi$(h/m) of the Bridge C period function (s07\_lz.tex:13) are cyclotomic. Quantum modular forms arising from false and indefinite theta series give WRT and Ẑ invariants: Zagier, "Quantum modular forms" (Clay Math Proc. 12, 2010); Lawrence--Zagier, Asian J Math 3 (1999) 93; Gukov--Pei--Putrov--Vafa, arXiv:1701.06567; Cheng--Chun--Ferrari--Gukov--Harrison, arXiv:1809.10148.
\item  \textbf{Next computation:} test whether $\psi$(h/m) at s = 7 satisfies a quantum-modular cocycle with real-analytic error. Search the plumbed or Seifert manifolds whose Ẑ uses an indefinite lattice of determinant -n\_s.
\end{itemize}
\end{itemize}

Also noted but weaker:
\begin{itemize}
\item  K3 lattices and Mathieu moonshine (Eguchi--Ooguri--Tachikawa, arXiv:1004.0956). The manuscript's K3 section addresses only lattice embedding; no elliptic genus appears.
\end{itemize}

\textbf{None of the three is an observable prediction.} They are mathematical correspondences that could become physics only through such a construction.

\end{document}